\documentclass{article}
\usepackage[utf8]{inputenc}
\usepackage{amsmath}
\usepackage{geometry}
\geometry{margin=2cm}
\begin{document}

\section*{Análise da Questão 157 — ENEM 2024 — Matemática}

\subsection*{Enunciado}
Uma doceira vende e entrega, em seu bairro, porções de 100 g de docinhos de aniversário. Atualmente, a taxa única de entrega é R$ 10,00, e o valor cobrado por uma porção é R$ 25,00. Por uma estratégia de vendas, a partir da próxima semana, a taxa única de entrega será R$ 15,00, e um novo valor será cobrado por uma porção, de maneira que o valor total a ser pago por um cliente na compra de 5 porções permaneça o mesmo.

A partir da próxima semana, qual será o novo valor cobrado, em real, por uma porção?

\subsection*{Solução Técnica}
\textbf{Resposta final:} R$24,00 \\
\textbf{Justificativa:} A resposta está correta e justificada pela resolução algébrica que mantém o total pago constante. \\
\textbf{Estratégia:} Igualar o custo total atual com o novo custo total, considerando o aumento da taxa fixa e ajustando o preço por porção. \\
\textbf{Passos da resolução:}
\begin{enumerate}
    \item Calcular o custo total atual:
    \begin{equation*}
    5 \times 25 + 10 = 135 \text{ reais}.
    \end{equation*}
    \item Expressar o custo total com o novo preço e nova taxa:
    \begin{equation*}
    5x + 15.
    \end{equation*}
    \item Igualar os totais:
    \begin{equation*}
    5x + 15 = 135.
    \end{equation*}
    \item Isolar $x$:
    \begin{equation*}
    5x = 120 \,\Rightarrow\, x = \frac{120}{5} = 24.
    \end{equation*}
    \item Concluir que o novo preço por porção é R$24,00.
\end{enumerate}

\subsection*{Explicação Didática}
\textbf{Explicação resumida:} O problema exige ajustar o preço por porção para que o custo total de 5 porções mais a taxa de entrega permaneça constante após o aumento da taxa. \\
\textbf{Passo a passo:}
\begin{enumerate}
    \item Calcular o custo total atual:
    \begin{equation*}
    5 \times 25 + 10 = 135.
    \end{equation*}
    \item Montar a equação para o novo custo total com o novo preço:
    \begin{equation*}
    5x + 15 = 135.
    \end{equation*}
    \item Isolar o termo com $x$:
    \begin{equation*}
    5x = 135 - 15 = 120.
    \end{equation*}
    \item Calcular o valor de $x$:
    \begin{equation*}
    x = \frac{120}{5} = 24.
    \end{equation*}
    \item Concluir que o novo preço por porção é R$24,00.
\end{enumerate}

\textbf{Erros comuns e dicas para o ensino foram descritos nas seções seguintes.}

\subsection*{Distratores}
\begin{itemize}
    \item \textbf{12,50}:
    \begin{itemize}
        \item Raciocínio provável: Dividir o total pago original apenas pelas porções, ignorando a taxa fixa.
        \item Por que é errado: Desconsidera a taxa fixa, portanto preço por porção incorreto.
    \end{itemize}
    \item \textbf{20,00}:
    \begin{itemize}
        \item Raciocínio provável: Subtrair diretamente o aumento da taxa do preço original proporcionalmente.
        \item Por que é errado: Não mantém o valor total constante conforme a equação correta.
    \end{itemize}
    \item \textbf{30,00}:
    \begin{itemize}
        \item Raciocínio provável: Aumentar o preço para cobrir a taxa extra, ignorando manter o gasto total fixo.
        \item Por que é errado: Contraria a condição de manter o custo total.
    \end{itemize}
    \item \textbf{37,50}:
    \begin{itemize}
        \item Raciocínio provável: Aplicar proporcionalmente a nova taxa de entrega ao preço, errado.
        \item Por que é errado: Não satisfaz condição de valor total inalterado.
    \end{itemize}
\end{itemize}

\subsection*{Perfil do Item}
Competências envolvidas:
\begin{itemize}
    \item Resolução de equações lineares (principal);
    \item Interpretação de problemas cotidianos;
    \item Operações com números reais.
\end{itemize}

Outras características:
\begin{itemize}
    \item Nível médio de dificuldade e demanda cognitiva;
    \item Não requer interpretação visual, mas requer modelagem e cálculos algébricos;
    \item Enfatiza tradução do contexto em equação e manipulação algébrica para solução.
\end{itemize}

\subsection*{Revisão de Consistência}
A análise técnica e a explicação didática estão consistentes entre si e com o enunciado. A solução correta (R$24,00) está bem justificada, e os distractores foram analisados adequadamente. O perfil do item está alinhado com as habilidades requeridas. Não foram identificados problemas de inconsistência ou erros relevantes.

\end{document}